\section{Effort to pulse width: the ESC model}
\label{sec:esc}

Each electronic speed controller (ESC) reads a pulse every 20\,ms (50\,Hz).
The pulse width $w$ in microseconds sets the thrust: $1500\,\mu$s is stopped,
longer is forward, shorter is reverse. The driver uses the linear map
\[
  w = \operatorname{clip}\!\left(1500 + 400\,t,\; 1100,\; 1900\right),
\]
so full effort $t=\pm1$ is $1900$ or $1100\,\mu$s. The pulse for thruster
position $p$ goes to physical output $c_p$ from the channel map.

The simulator inverts the same model. A pulse $w$ on an output becomes a motor
command
\[
  m = m_\text{max}\,\frac{w - 1500}{400},\qquad w\in[1100,1900],
\]
optionally with a deadband $|w-1500|\le d \Rightarrow m = 0$, which real ESCs
have. Because both sides share the model, a characterization done in simulation
exercises the same chain as on the boat.

\paragraph{Why republish.} The simulator treats a silent pulse topic as a lost
link and goes neutral after a timeout, as the real ESCs do without a signal. The
driver therefore repeats the last pulses at 20\,Hz even when they do not change.
\file{mhseals\_hardware/thruster\_mixer.py (effort\_to\_pulse), outputs.py; crane\_ml EscPwm.cs}
